the case of an ideal in a ring.

\begin{definition}
\label{definition-regular-ideal}
Let $R$ be a ring and let $I \subset R$ be an ideal.
\begin{enumerate}
\item We say $I$ is a {\it regular ideal} if for every
$\mathfrak p \in V(I)$ there exists a $g \in R$, $g \not \in \mathfrak p$
and a regular sequence $f_1, \ldots, f_r \in R_g$ such that $I_g$
is generated by $f_1, \ldots, f_r$.
\item We say $I$ is a {\it Koszul-regular ideal} if for every
$\mathfrak p \in V(I)$ there exists a $g \in R$, $g \not \in \mathfrak p$
and a Koszul-regular sequence $f_1, \ldots, f_r \in R_g$ such that $I_g$
is generated by $f_1, \ldots, f_r$.
\item We say $I$ is a {\it $H_1$-regular ideal} if for every
$\mathfrak p \in V(I)$ there exists a $g \in R$, $g \not \in \mathfrak p$
and an $H_1$-regular sequence $f_1, \ldots, f_r \in R_g$ such that $I_g$
is generated by $f_1, \ldots, f_r$.
\item We say $I$ is a {\it quasi-regular ideal} if for every
$\mathfrak p \in V(I)$ there exists a $g \in R$, $g \not \in \mathfrak p$
and a quasi-regular sequence $f_1, \ldots, f_r \in R_g$ such that $I_g$
is generated by $f_1, \ldots, f_r$.
\end{enumerate}
\end{definition}

\noindent
It is clear that given $I \subset R$ we have the implications
\begin{align*}
I\text{ is a regular ideal}
& \Rightarrow
I\text{ is a Koszul-regular ideal} \\
& \Rightarrow
I\text{ is a }H_1\text{-regular ideal} \\
& \Rightarrow
I\text{ is a quasi-regular ideal}
\end{align*}
see Lemmas \ref{lemma-regular-koszul-regular},
\ref{lemma-koszul-regular-H1-regular}, and
\ref{lemma-H1-regular-quasi-regular}. Such an ideal is always finitely
generated.

\begin{lemma}
\label{lemma-quasi-regular-ideal-finite}
A quasi-regular ideal is finitely generated.
\end{lemma}

\begin{proof}
Let $I \subset R$ be a quasi-regular ideal. Since $V(I)$ is quasi-compact,
there exist $g_1, \ldots, g_m \in R$ such that
$V(I) \subset D(g_1) \cup \ldots \cup D(g_m)$
and such that $I_{g_j}$ is generated by a quasi-regular sequence
$g_{j1}, \ldots, g_{jr_j} \in R_{g_j}$. Write $g_{ji} = g'_{ji}/g_j^{e_{ij}}$
for some $g'_{ij} \in I$. Write $1 + x = \sum g_j h_j$
for some $x \in I$ which is possible as
$V(I) \subset D(g_1) \cup \ldots \cup D(g_m)$.
Note that $\Spec(R) = D(g_1) \cup \ldots \cup D(g_m) \bigcup D(x)$
Then $I$ is generated by the elements $g'_{ij}$ and $x$ as
these generate on each of the pieces of the cover, see
Algebra, Lemma \ref{algebra-lemma-cover}.
\end{proof}

\begin{lemma}
\label{lemma-quasi-regular-ideal-finite-projective}
Let $I \subset R$ be a quasi-regular ideal of a ring.
Then $I/I^2$ is a finite projective $R/I$-module.
\end{lemma}

\begin{proof}
This follows from Algebra, Lemma \ref{algebra-lemma-finite-projective}
and the definitions.
\end{proof}

\noindent
We prove flat descent for Koszul-regular, $H_1$-regular, quasi-regular
ideals.

\begin{lemma}
\label{lemma-flat-descent-regular-ideal}
Let $A \to B$ be a faithfully flat ring map. Let $I \subset A$ be an ideal.
If $IB$ is a Koszul-regular
(resp.\ $H_1$-regular, resp.\ quasi-regular) ideal in $B$, then
$I$ is a Koszul-regular (resp.\ $H_1$-regular, resp.\ quasi-regular)
ideal in $A$.
\end{lemma}

\begin{proof}
We fix the prime $\mathfrak p \supset I$ throughout the proof.
Assume $IB$ is quasi-regular. By
Lemma \ref{lemma-quasi-regular-ideal-finite}
$IB$ is a finite module, hence $I$ is a finite $A$-module by
Algebra, Lemma \ref{algebra-lemma-descend-properties-modules}.
As $A \to B$ is flat we see that
$$
I/I^2 \otimes_{A/I} B/IB = I/I^2 \otimes_A B = IB/(IB)^2.
$$
As $IB$ is quasi-regular, the $B/IB$-module $IB/(IB)^2$ is finite
locally free. Hence $I/I^2$ is finite projective, see
Algebra, Proposition \ref{algebra-proposition-ffdescent-finite-projectivity}.
In particular, after replacing $A$ by $A_f$ for some
$f \in A$, $f \not \in \mathfrak p$ we may assume that $I/I^2$ is free of
rank $r$. Pick $f_1, \ldots, f_r \in I$ which
give a basis of $I/I^2$. By Nakayama's lemma (see
Algebra, Lemma \ref{algebra-lemma-NAK})
we see that, after another replacement $A \leadsto A_f$ as above,
